%!TEX root = tfg.tex
\chapter{Introducción}

\begin{quotation}[Poet]{Dante Alighieri (1265--1321)}
Midway upon the journey of our life I found myself within a forest dark, for the straightforward pathway had been lost.
\end{quotation}

\begin{abstract}
Este capítulo presenta una visión general de \textit{Via Inferni} y sitúa el proyecto
dentro del marco de un Trabajo de Fin de Grado de Ingeniería del Software. Se resume la
propuesta del videojuego, la motivación que justifica su desarrollo y la organización
general de la memoria.
\end{abstract}

\section{Presentación del proyecto}
\textit{Via Inferni} es un videojuego \textit{roguelike} 2D con generación procedural
de niveles inspirado en el Infierno descrito por Dante Alighieri en la
\textit{Divina Comedia}. La propuesta combina mecánicas habituales del género, como la
exploración de salas, la progresión por intentos y el combate en tiempo real, con una
ambientación literaria que organiza el avance del jugador a través de diferentes círculos
temáticos.

Desde el punto de vista académico, este TFG se plantea como un ejercicio integral de
ingeniería del software. No solo se aborda la implementación del producto, sino también
su planificación, la gestión del alcance, la organización del trabajo, el control del
código y la documentación del proceso seguido hasta obtener una versión jugable del
sistema. El proyecto ha sido desarrollado por dos estudiantes de Ingeniería del Software,
por lo que la memoria busca reflejar tanto el resultado final como las decisiones
técnicas y de gestión tomadas durante el trabajo.

\section{Motivación}
La motivación principal del proyecto es aplicar los conocimientos adquiridos durante el
grado en un producto software completo, con restricciones reales de tiempo, alcance y
coordinación. Un videojuego resulta especialmente adecuado para este propósito porque
obliga a integrar áreas diversas: arquitectura de software, interacción con el usuario,
gestión de estados, persistencia, diseño de niveles, pruebas y control de versiones.

También existe una motivación creativa. La estructura de los círculos del Infierno ofrece
una base clara para organizar niveles, enemigos y progresión de dificultad, mientras que
el formato \textit{roguelike} permite construir partidas breves, variables y con alto
grado de rejugabilidad. Esta combinación permite trabajar sobre un producto técnicamente
exigente sin perder una identidad temática reconocible.

Por último, el proyecto responde al interés de sus autores por el desarrollo de
videojuegos y por la aplicación práctica de una metodología de ingeniería en un entorno
cercano al desarrollo independiente. La intención no es producir un videojuego comercial
completo, sino un prototipo funcional y documentado que permita evaluar el proceso seguido
y el grado de cumplimiento de los objetivos definidos.

\section{Estructura de la memoria}
La memoria se ha organizado en cinco grandes bloques. El primero, \textbf{Inicio},
presenta el proyecto, su contexto y sus objetivos. El segundo,
\textbf{Planificación}, desarrolla los distintos planes de gestión y la
organización seguida durante el desarrollo del proyecto. El tercer bloque,
\textbf{Ejecución, seguimiento y control}, documenta el proceso de desarrollo
iterativo del videojuego. El cuarto, \textbf{Cierre}, recoge los resultados
obtenidos, las conclusiones y las instrucciones de uso e instalación. Por
último, los \textbf{Apéndices} reúnen la documentación complementaria necesaria
para completar el trabajo.
